% title: Weighted Binomial Sum
% short-title: Weighted Binomial Sum
% abstract: We formalize in Lean 4 a classical weighted binomial identity: the sum of the binomial coefficients in the $(n+1)$st row of Pascal's triangle, weighted by the index, equals $(n+1)2^n$.
% subjclass: 05A10
% keywords: binomial coefficient, weighted sum, Pascal triangle, Lean 4
% author: Mario Hernández M.

\subsection{Weighted Binomial Sum}

The formal result in this section is
\lean{Biblioteca.Demonstrations.weighted_binomial_sum}. It states that for
every natural number $n$,
\[
  \sum_{i=0}^{n+1} i \binom{n+1}{i} = (n+1) 2^n.
\]

\begin{theorem}
For every $n \in \mathbb{N}$,
\[
  \sum_{i=0}^{n+1} i \binom{n+1}{i} = (n+1) 2^n.
\]
\end{theorem}

\begin{proof}
This is the classical first moment identity for the binomial coefficients. A
standard combinatorial argument counts the number of ways to choose a nonempty
subset of an $(n+1)$-element set together with one distinguished element of
that subset. If one first chooses the subset size $i$, then there are
$i\binom{n+1}{i}$ possibilities. Summing over all sizes gives the left-hand
side.

On the other hand, one may first choose the distinguished element, which gives
$n+1$ choices, and then choose any subset of the remaining $n$ elements to join
it. This contributes $(n+1)2^n$ possibilities, proving the identity.

The Lean formalization uses the library theorem
\texttt{Nat.sum\_range\_mul\_choose} and rewrites it into the shifted form above.
\end{proof}
